\chapter*{Abstract}
\addcontentsline{toc}{chapter}{Abstract}
\markboth{Abstract}{Abstract}

This Bachelor's Thesis studies the automatic analysis of electrocardiograms (ECGs) represented as images. It focuses on suspected Brugada syndrome, a rare disease for which large labelled patient datasets are difficult to assemble. The aim is to examine how a few examples can help identify relevant ECG patterns.

The two strategies compared are supervised training of a ResNet18 convolutional neural network (CNN) and in-context learning (ICL) with the Gemma~4 vision-language model, which receives labelled examples in the query without weight updates. Both analyse an image of one heartbeat from lead V1, one of the ECG traces recorded from the chest. Experiments use between 2 and 32 examples, with 100 QRS/ST simulator cases and ECGs from 317 patients at Hospital Universitario Central de Asturias (HUCA).

First, examples and test images from the simulator assess recognition of three predefined waveform features. With 32 examples, ResNet18 reaches a balanced accuracy (BA) of 0.885 and Gemma~4 reaches 0.904 with equally represented classes. Performance falls when simulator-based learning is applied to clinical ECGs. The CNN achieves a maximum BA of 0.551 on HUCA, while Gemma~4 reaches 0.514 with simulator demonstrations and tends to predict the negative class.

Finally, HUCA examples are used to classify other patients from the same dataset. A second CNN trained directly on these data reaches a BA of 0.546 with 16 examples. It remains near chance and slightly outperforms Gemma~4 with 2 and 4 examples. Gemma~4 performs better from 8 onwards, reaching a BA of 0.761 with 32 balanced HUCA demonstrations, 0.741 sensitivity and 0.781 specificity.

Without examples, Gemma~4 identifies no positives in either dataset. HUCA demonstrations help it classify other patients from the same dataset, whereas simulator examples do not adequately transfer this learning to real ECGs. These results support investigating triage assistance to prioritise review and collect data for models with better performance or explainability. Such assistance would require clinical validation and would not replace a full medical assessment.

\textbf{Keywords:} electrocardiogram as image, Brugada syndrome, few-shot learning, convolutional neural networks, vision-language models, in-context learning, domain adaptation.
